% Chapter 9: Universal Learning and the Trichotomy
% universal_learning = Profile C: what is the optimal learning RATE over all distributions?
% universal_trichotomy = Profile C: the modern capstone -- three regimes,
%   classified by Littlestone trees (from online learning!) applied to the statistical setting.
% Centerpiece: the SURPRISE that the Littlestone dimension controls statistical rates.

\chapter{Universal Learning and the Trichotomy}
\label{ch:universal}

The Fundamental Theorem of \Cref{ch:pac} answers a binary question: \emph{is}
$\mathcal{H}$ learnable?  The answer is yes if and only if
$\VC(\mathcal{H}) < \infty$.  But the binary answer hides a quantitative
question.  Among all learnable classes, some are learned faster than others.
How fast?

PAC learning theory phrases this quantitatively through the sample complexity
$m(\varepsilon, \delta)$: how many samples suffice for accuracy $\varepsilon$
with confidence $1 - \delta$?  But the answer, $\Theta(d/\varepsilon)$ in
the realizable case, is a function of the \emph{accuracy parameter}.  It
describes how accuracy improves as data grows, with $d$ as a constant.  It
does not describe the \emph{rate} at which the risk itself decreases as a
function of the sample size $n$.

This chapter asks the rate question directly: given $n$ samples from an
unknown distribution $D$, how fast can the risk $\risk_D(A(S))$ be driven to
zero, \emph{uniformly over all distributions $D$ for which learning is
possible}?  The answer is the trichotomy theorem, which classifies every
hypothesis class into exactly one of three rate regimes.  The classification
involves a combinatorial object from an unexpected source.


% ================================================================
% SECTION 1: THE QUESTION
% ================================================================

\section{Beyond PAC Sample Complexity}
\label{sec:universal-question}

Fix a hypothesis class $\mathcal{H}$ with $\VC(\mathcal{H}) < \infty$, so
that $\mathcal{H}$ is PAC learnable.  The PAC framework guarantees a learner
$A$ and a sample complexity $m(\varepsilon, \delta)$ such that, for any
distribution $D$ and any $\varepsilon, \delta > 0$, drawing
$m \geq m(\varepsilon, \delta)$ samples ensures
$\Pr[\risk_D(A(S)) > \varepsilon] \leq \delta$.

Now invert the perspective.  Instead of asking ``how many samples for
accuracy $\varepsilon$?'', ask: ``given $n$ samples, what is the best
accuracy achievable?''  More precisely, define the \emph{learning rate} of a
class $\mathcal{H}$ as follows.

\begin{defn}[Universal Learning Rate]
\label{def:universal-rate}
A function $R \colon \N \to [0,1]$ is a \emph{universal learning rate} for
$\mathcal{H}$ if there exists a learner $A$ such that, for every realizable
distribution $D$ (i.e., $\risk_D(h^*) = 0$ for some $h^* \in \mathcal{H}$)
and every $n \in \N$:
\[
  \E_{S \sim D^n}\!\bigl[\risk_D(A(S))\bigr] \leq R(n).
\]
The \emph{optimal universal rate} is the fastest-decaying $R$ achievable by
any learner.\formal{FLT_Proofs.Criterion.Extended.UniversalLearnable}
\end{defn}

Three features distinguish this from the PAC formulation:

\begin{enumerate}[leftmargin=*,itemsep=3pt]
\item \textbf{No fixed $\varepsilon$.}  The rate $R(n)$ describes how the
  expected risk shrinks with $n$, rather than asking for a threshold sample
  size at a given $\varepsilon$.

\item \textbf{Uniformity over distributions.}  The same learner $A$ and the
  same rate $R(n)$ must work for \emph{every} realizable distribution $D$.
  The learner cannot be tuned to a specific distribution.

\item \textbf{Expected risk, not high-probability.}  We use
  $\E[\risk_D(A(S))]$ rather than a $(\varepsilon,\delta)$ guarantee.  This
  is a convenience that simplifies the rate classification without changing
  the qualitative picture.
\end{enumerate}

For a class with $\VC(\mathcal{H}) = d < \infty$, the PAC upper bound gives
$R(n) = O(d/n)$: a $\Theta(1/n)$ rate.  But is $1/n$ always achievable?
Can some classes be learned faster, exponentially fast?  And are there
learnable classes where no uniform $1/n$ rate is possible, so that the rate
degrades depending on the distribution?

All three phenomena occur.  The trichotomy theorem says these are the
\emph{only} three possibilities.


% ================================================================
% SECTION 2: THE TRICHOTOMY THEOREM
% ================================================================

\section{The Trichotomy Theorem}
\label{sec:trichotomy}

The classification involves a combinatorial object that generalizes the
Littlestone tree of \Cref{ch:online}.  Recall that a mistake tree
(\Cref{def:mistake-tree}) is a complete binary tree whose internal nodes are
labeled with instances $x \in X$, such that every root-to-leaf path is
realized by some $h \in \mathcal{H}$.  We now define a variant with a
weaker consistency requirement.

\begin{defn}[VCL Tree]
\label{def:vcl-tree}
A \emph{VCL tree} (Vapnik--Chervonenkis--Littlestone tree) for
$\mathcal{H}$ over $X$ is a complete binary tree $T$ in which:
\begin{itemize}[leftmargin=*,itemsep=2pt]
\item Each internal node $v$ is labeled with an instance $x_v \in X$.
\item The left child corresponds to label $0$, the right to label $1$.
\item For every root-to-leaf path $\pi = (v_1, y_1), \ldots, (v_d, y_d)$,
  there exists a \emph{set} $\{x_{v_1}, \ldots, x_{v_d}\}$ that is
  \emph{shattered} by $\mathcal{H}$, not merely a single hypothesis
  consistent with $\pi$, but the full shattering condition: for every
  labeling $b \in \{0,1\}^d$ of these $d$ points, some $h \in \mathcal{H}$
  realizes $b$.
\end{itemize}
The tree is \emph{infinite} if it has unbounded depth.
\end{defn}

\begin{remark}[VCL trees vs.\ Littlestone trees]
\label{rmk:vcl-vs-littlestone}
Every Littlestone tree is a VCL tree (shattering is a stronger condition
than path-consistency, but the shattering requirement in the VCL definition
applies to each path individually, and a Littlestone tree already provides a
consistent hypothesis for every path, since the tree is complete, all
$2^d$ paths are realized, which implies shattering of the instances along
any single path).  The converse fails: a VCL tree requires shattering along
each path, which is a weaker global condition than the Littlestone tree's
requirement that \emph{every} root-to-leaf path has a \emph{single}
consistent hypothesis.  The distinction matters: infinite VCL trees can
exist even when $\Ldim(\mathcal{H}) < \infty$.  The combined VC--Littlestone measure underlying these trees is typed in the development as a pair of ordinal-valued dimensions.\formal{FLT_Proofs.Complexity.Ordinal.OrdinalLittlestoneDim}
\end{remark}

\begin{thm}[The Trichotomy Theorem {\cite{BousquetHannekeMoranVanHandelYehudayoff2021}}]
\label{thm:trichotomy}
Let $\mathcal{H}$ be a hypothesis class with $|\mathcal{H}| \geq 3$.
Exactly one of the following three cases holds:
\begin{enumerate}[label=\textnormal{(T\arabic*)}]
\item \label{T:exp} \textbf{Exponential rate.}  $\mathcal{H}$ has no
  infinite Littlestone tree (equivalently, $\Ldim(\mathcal{H}) < \infty$).
  Then the optimal universal learning rate is \emph{exponential}:
  \[
    R(n) = \Theta\!\left(e^{-n}\right).
  \]

\item \label{T:linear} \textbf{Linear rate.}  $\mathcal{H}$ has an infinite
  Littlestone tree but no infinite VCL tree.  Then the optimal universal
  learning rate is \emph{linear}:
  \[
    R(n) = \Theta\!\left(\frac{1}{n}\right).
  \]

\item \label{T:arb} \textbf{Arbitrarily slow rates.}  $\mathcal{H}$ has an
  infinite VCL tree.  Then no uniform rate exists: for every function
  $R(n) \to 0$ (no matter how slowly), every learner $A$ fails to achieve
  rate $R(n)$ on some realizable distribution $D$.
\end{enumerate}
\end{thm}

\begin{remark}[Formalization status]
\label{rmk:trichotomy-formal}
The trichotomy's three \emph{regimes} are typed in the companion development (the universal-rate criterion (\cref{def:universal-rate}) and the ordinal VC and Littlestone dimensions), and the boundary case ``$\Ldim < \infty$ iff online learnable'' is verified (\cref{thm:littlestone-characterization}).  The trichotomy theorem itself is \emph{not} yet machine-checked: the middle branch (infinite Littlestone, finite VC, hence linear rate) rests on the Bousquet--Hanneke--Moran--Ramanan doubling construction, which the development marks as pending, and the universal-to-PAC direction routes through a boosting lemma that is currently incomplete.  Closing these is \cref{op:trichotomy-formal}; we present the proof as a sketch throughout, following~\cite{BousquetHannekeMoranVanHandelYehudayoff2021}.
\end{remark}

The three cases are exhaustive and mutually exclusive.  Every learnable class
(i.e., every class with $\VC(\mathcal{H}) < \infty$ and
$|\mathcal{H}| \geq 3$) falls into exactly one regime.  The boundaries are
sharp: there is no class whose optimal rate is, say, $1/\sqrt{n}$ or
$1/n^2$.

\begin{figure}[ht]
\centering
\begin{tikzpicture}[
    decision/.style={draw, thick, diamond, aspect=2.2, inner sep=2pt, font=\small,
      fill=defblue!8, draw=defblue!70!black, align=center, minimum width=3cm},
    outcome/.style={draw, thick, densely dashed, rounded corners, minimum width=2.7cm,
      minimum height=1.25cm, font=\small, align=center},
    sedge/.style={-{Stealth[length=2.5mm]}, thick},
    dedge/.style={-{Stealth[length=2.5mm]}, thick, densely dashed},
    yeslabel/.style={font=\scriptsize\bfseries, text=sepgreen!60!black},
    nolabel/.style={font=\scriptsize\bfseries, text=thmred},
    hdl/.style={font=\tiny, text=defblue!55!black, align=center},
]
% --- decisions: verified dimensional thresholds (solid) ---
\node[decision] (d1) at (0,0) {Infinite\\Littlestone tree?};
\node[hdl, below=1pt of d1] {\leanname{OrdinalLittlestoneDim}};
\node[decision] (d2) at (3.6,-3.1) {Infinite\\VCL tree?};
\node[hdl, below=1pt of d2] {\cref{def:vcl-tree}, \leanname{VCLTree}};
% --- rate outcomes: typed regimes, but the rate THEOREMS are pending -> dashed boxes ---
\node[outcome, fill=sepgreen!12, draw=sepgreen!60!black] (exp) at (-3.5,-6.1) {%
  \textbf{Exponential}\\$R(n)=\Theta(e^{-n})$\\{\scriptsize $\Ldim<\infty$}};
\node[outcome, fill=defblue!12, draw=defblue!60!black] (lin) at (1.5,-6.1) {%
  \textbf{Linear}\\$R(n)=\Theta(1/n)$\\{\scriptsize $\Ldim=\infty$, no $\infty$-VCL}};
\node[outcome, fill=thmred!12, draw=thmred!60!black] (slow) at (6.3,-6.1) {%
  \textbf{Arbitrarily slow}\\$\forall\,R(n)\to0$: fails\\{\scriptsize $\infty$-VCL exists}};
% --- edges: D1 solid (verified boundary), D2 dashed (rate consequence is literature) ---
\draw[sedge] (d1) -- node[nolabel, above left, pos=0.35]{No}
   node[hdl, below, pos=0.62]{\leanname{littlestone_characterization}} (exp);
\draw[sedge] (d1) -- node[yeslabel, above right, pos=0.4]{Yes} (d2);
\draw[dedge] (d2) -- node[nolabel, above left, pos=0.4]{No} (lin);
\draw[dedge] (d2) -- node[yeslabel, above right, pos=0.4]{Yes} (slow);
% --- the inclusion ladder + the formalization-status annotation ---
\node[hdl, text=black!55, anchor=west] at (4.55,-1.15)
   {$\{\Ldim{<}\infty\}\subsetneq\{$no $\infty$-VCL$\}\subsetneq\{$all$\}$};
\node[font=\scriptsize\itshape, text=thmred!75!black, align=center, anchor=north, text width=11cm] at (1.4,-7.5)
   {Solid: the decisions are verified thresholds, and \leanname{littlestone_characterization} is the one verified rate boundary. Dashed: the $\Theta$-rates and the trichotomy are \cite{BousquetHannekeMoranVanHandelYehudayoff2021}, typed but not yet machine-checked (\cref{op:trichotomy-formal}).};
\end{tikzpicture}
\caption[The universal-rate trichotomy: two tree-conditions, three regimes]{The trichotomy decision tree.
Two combinatorial questions partition every hypothesis class with $|\mathcal{H}|\geq 3$ into exactly three
universal learning-rate regimes. The first question, does $\mathcal{H}$ admit an infinite Littlestone
tree, is a verified dimensional threshold: by the Littlestone characterization
(\cref{thm:littlestone-characterization}, \leanname{littlestone_characterization}) a finite Littlestone
dimension is exactly online learnability, and this is the boundary of the exponential regime
$R(n)=\Theta(e^{-n})$. The second question, does $\mathcal{H}$ admit an infinite VCL tree, uses the VCL
tree of \cref{def:vcl-tree}, which generalizes the Littlestone tree by requiring that the instances along
each root-to-leaf path be shattered rather than merely consistent with one hypothesis; its finiteness
separates the linear regime $R(n)=\Theta(1/n)$ from the arbitrarily-slow regime. The two decision nodes
are drawn solid because the dimensions and the VCL tree are typed in the companion development; the three
rate values and the ``exactly one of three'' claim are drawn dashed because they are the result of
Bousquet--Hanneke--Moran--van Handel--Yehudayoff~\cite{BousquetHannekeMoranVanHandelYehudayoff2021}, cited
here and not yet machine-checked (\cref{op:trichotomy-formal}). There is no fourth regime and no
intermediate rate.}
\label{fig:trichotomy-tree}
\end{figure}

\subsection{An Online Concept in a Statistical Theorem}
\label{sec:surprise}

The trichotomy theorem is the structural capstone of this book.  To
appreciate what it says, consider the expectations each tradition brought
to the table.

The PAC tradition (\Cref{ch:pac}) characterized learnability through the VC
dimension and established the sample complexity
$\Theta(d/\varepsilon)$.  The natural expectation: the VC dimension should
also determine the learning rate.  It does not.
Two classes with the same VC dimension can have different optimal
rates: one exponential, one $1/n$.  So what draws the boundary?

The online learning tradition (\Cref{ch:online}) introduced the Littlestone
dimension to characterize mistake-bounded learning against an adversary.
A worst-case combinatorial quantity, built for a worst-case adversarial
game.  No connection to statistical convergence rates.  None expected.

And yet.  The Littlestone dimension, this adversarial, combinatorial
object, is exactly what governs the boundary between exponential and
polynomial convergence in the i.i.d.\ setting.

Neither tradition predicted this.  The trichotomy does not merely classify
rates; it reveals that the deepest structural boundary in statistical
learning theory is drawn by a quantity from a different paradigm entirely.

\subsection{Examples}
\label{sec:trichotomy-examples}

\begin{example}[The three regimes, witnessed]
\label{ex:three-regimes}
\begin{enumerate}[label=(\alph*),leftmargin=*,itemsep=4pt]
\item \textbf{Exponential: finite classes.}  Any finite class
  $\mathcal{H}$ with $|\mathcal{H}| = k$ has
  $\Ldim(\mathcal{H}) \leq \lfloor\log_2 k\rfloor < \infty$.  No infinite
  Littlestone tree exists, so $\mathcal{H}$ falls into regime~\ref{T:exp}.
  Concretely, the Halving algorithm achieves risk $\leq k \cdot 2^{-n}$
  after $n$ rounds.

\item \textbf{Exponential: halfspaces.}  The class of halfspaces in
  $\R^d$ has $\Ldim = d < \infty$ (\Cref{ex:ldim-examples}).  This class
  is learned at an exponential rate, despite having infinite cardinality.

\item \textbf{Linear: thresholds on $\R$.}  The class
  $\mathcal{H}_{\mathrm{thr}} = \{x \mapsto \ind[x \geq \theta] :
  \theta \in \R\}$ has $\VC = 1$ but $\Ldim = \infty$
  (\Cref{prop:pac-online-sep}).  It has no infinite VCL tree (the
  one-dimensional ordering prevents shattering along arbitrarily long
  adaptive paths).  The optimal rate is $\Theta(1/n)$.

\item \textbf{Arbitrarily slow: unions of intervals.}  For each $k$, let
  $\mathcal{H}_k$ be the class of unions of at most $k$ intervals on $\R$.
  The class $\mathcal{H}_\infty = \bigcup_k \mathcal{H}_k$ (unions of
  finitely many intervals, with no bound on the number) has
  $\VC < \infty$ but admits infinite VCL trees.  No uniform rate $R(n) \to 0$
  is achievable.
\end{enumerate}
\end{example}


% ================================================================
% SECTION 3: PROOF SKETCH
% ================================================================

\section{Proof Architecture}
\label{sec:proof-sketch}

The full proof of \Cref{thm:trichotomy} occupies much of
\cite{BousquetHannekeMoranVanHandelYehudayoff2021}.  We sketch the key ideas
for each regime, emphasizing why the boundaries fall where they do.


\subsection{The Exponential Case: Finite Littlestone Dimension}
\label{sec:exp-case}

Suppose $\Ldim(\mathcal{H}) = d < \infty$.  The Standard Optimal Algorithm
($\SOA$) from \Cref{ch:online} was designed for adversarial online learning,
but its structure can be adapted to the statistical setting.

The key insight is that a learner with access to i.i.d.\ samples can
simulate a version-space strategy.  Consider the following approach: given a
sample $S$ of size $n$, maintain the version space
$V = \{h \in \mathcal{H} : h \text{ consistent with } S\}$.  Return a
hypothesis from $V$ using an $\SOA$-derived selection rule.

When $\Ldim(\mathcal{H}) = d$, the version space after seeing $n$ i.i.d.\
samples has the following property: the probability (over $D$) that there
exists a hypothesis $h \in V$ with $\risk_D(h) > \varepsilon$ decreases
exponentially in $n$.  This is because each sample point, with probability
at least $\varepsilon$ under $D$, eliminates all hypotheses in $V$ that
disagree with the true label at that point.  After $n$ samples, the survival
probability of any ``bad'' hypothesis (one with risk $> \varepsilon$) is at
most $(1 - \varepsilon)^n \leq e^{-\varepsilon n}$.  The finite Littlestone
dimension controls the effective complexity of $V$ through the mistake tree
structure, ensuring that a union bound over the relevant version space
subsets remains finite.

The result:
\[
  \E[\risk_D(A(S))] \leq C \cdot 2^d \cdot e^{-n/C'}
\]
for constants $C, C'$ depending on $d$.  The rate is exponential in $n$.

The matching lower bound shows that exponential convergence is impossible
when $\Ldim(\mathcal{H}) = \infty$: an infinite Littlestone tree provides an
adversary-like construction (within the distributional framework) that
forces the learner's expected risk to decay no faster than $\Omega(1/n)$.


\subsection{The Linear Case: Infinite Littlestone, Finite VCL}
\label{sec:linear-case}

Suppose $\Ldim(\mathcal{H}) = \infty$ but $\mathcal{H}$ has no infinite VCL
tree.  The exponential strategy fails because the $\SOA$-based approach
requires finite Littlestone dimension.  But a different algorithm, the
\emph{one-inclusion predictor}, achieves the $1/n$ rate.

\begin{defn}[One-Inclusion Predictor (Informal)]
\label{def:oip}
Given a sample $S = \{(x_1, y_1), \ldots, (x_n, y_n)\}$ and a new point
$x$, form the multiset $\{x_1, \ldots, x_n, x\}$.  Consider the one-inclusion
graph $G$: the graph whose vertices are all labelings of this multiset
consistent with some $h \in \mathcal{H}$, with edges connecting labelings
that differ on exactly one point.  Predict the label $\hat{y}$ that
minimizes the maximum degree of the vertex in $G$ corresponding to the
completed labeling.
\end{defn}

Haussler, Littlestone, and Warmuth~\cite{Haussler1988} introduced the
one-inclusion graph technique.  Bousquet et
al.~\cite{BousquetHannekeMoranVanHandelYehudayoff2021} showed that when
$\mathcal{H}$ has no infinite VCL tree, the one-inclusion predictor achieves
\[
  \E[\risk_D(A(S))] \leq \frac{C}{n}
\]
for a constant $C$ depending on $\mathcal{H}$.  The absence of infinite VCL
trees is precisely the combinatorial condition needed for the one-inclusion
graph to have bounded degree growth, which in turn controls the expected
risk.

The matching lower bound is the PAC lower bound: any class with
$\VC(\mathcal{H}) \geq 1$ requires $\Omega(1/n)$ expected risk, since
distinguishing hypotheses that differ on a $1/n$-fraction of the domain
requires $\Omega(n)$ samples.


\subsection{The Slow Case: Infinite VCL Trees}
\label{sec:slow-case}

Suppose $\mathcal{H}$ has an infinite VCL tree.  The theorem claims that no
uniform rate $R(n) \to 0$ is achievable: for every learner $A$ and every
function $R(n) \to 0$, there exists a realizable distribution $D$ such that
$\E[\risk_D(A(S_n))] \geq R(n)$ for infinitely many $n$.

The proof constructs, for any candidate learner $A$ and any target rate
$R(n)$, a distribution that defeats $A$ at the desired rate.  The
construction uses the infinite VCL tree as a source of combinatorial
richness.

Descend the VCL tree adaptively, choosing at each node the branch that is
harder for $A$.  Because the tree is infinite and each path corresponds to a
shattered set, the adversary can build a distribution at any depth $d$ whose
effective VC dimension on the relevant support is $d$.  By choosing $d$
large enough (as a function of $n$), the adversary ensures that the PAC
lower bound $\Omega(d/n)$ exceeds $R(n)$.  The infinite VCL tree provides
arbitrarily large effective VC dimension along adaptive paths, preventing
any uniform rate from holding.

This is the sense in which infinite VCL trees are an obstruction: they
provide an inexhaustible reservoir of complexity that any learner eventually
confronts.


% ================================================================
% SECTION 4: THE CROSS-PARADIGM MAP
% ================================================================

\section{The Cross-Paradigm Map}
\label{sec:cross-paradigm}

The trichotomy theorem connects three chapters of this book through a single
combinatorial object.

\begin{figure}[ht]
\centering
\begin{tikzpicture}[
    regime/.style={draw, thick, ellipse, minimum width=3.6cm,
      minimum height=2.2cm, font=\small, align=center},
    example/.style={font=\scriptsize, align=center, text=notegray},
    label/.style={font=\small\bfseries, align=center},
    connection/.style={thick, dashed, notegray},
]
% Three regimes as overlapping regions
\node[regime, fill=sepgreen!10] (exp) at (-3.5, 0) {};
\node[label] at (-3.5, 0.4) {$e^{-n}$};
\node[example] at (-3.5, -0.3) {finite classes\\halfspaces\\conjunctions};

\node[regime, fill=defblue!10] (lin) at (0, 0) {};
\node[label] at (0, 0.4) {$1/n$};
\node[example] at (0, -0.3) {thresholds on $\R$\\intervals on $\R$};

\node[regime, fill=thmred!10] (slow) at (3.5, 0) {};
\node[label] at (3.5, 0.4) {arb.\ slow};
\node[example] at (3.5, -0.3) {$\bigcup_k \mathcal{H}_k$\\unbounded\\unions};

% Boundaries
\draw[very thick, sepgreen!70!black, {Stealth[length=2mm]}-{Stealth[length=2mm]}]
  (-1.6, -1.8) -- (-1.6, 1.8)
  node[above, font=\scriptsize, align=center, text=sepgreen!70!black]
  {$\Ldim < \infty$\\vs.\ $\Ldim = \infty$};

\draw[very thick, thmred!70!black, {Stealth[length=2mm]}-{Stealth[length=2mm]}]
  (1.7, -1.8) -- (1.7, 1.8)
  node[above, font=\scriptsize, align=center, text=thmred!70!black]
  {no $\infty$-VCL\\vs.\ $\infty$-VCL};

% Chapter connections
\node[font=\scriptsize\itshape, text=defblue, align=center] at (-3.5, -2.6)
  {Ch.~6: Online\\(Littlestone dim)};
\node[font=\scriptsize\itshape, text=defblue, align=center] at (0, -2.6)
  {Ch.~5: PAC\\(VC dim)};
\node[font=\scriptsize\itshape, text=defblue, align=center] at (3.5, -2.6)
  {Ch.~9: Universal\\(VCL trees)};

\draw[connection, -{Stealth[length=1.5mm]}] (-3.5, -2.2) -- (-1.8, -1.6);
\draw[connection, -{Stealth[length=1.5mm]}] (0, -2.2) -- (0, -1.2);
\draw[connection, -{Stealth[length=1.5mm]}] (3.5, -2.2) -- (1.9, -1.6);
\end{tikzpicture}
\caption{The three rate regimes and their governing combinatorial conditions.
The left boundary is drawn by the Littlestone dimension (from online
learning, \Cref{ch:online}); the right boundary by VCL trees (a hybrid of
VC shattering and Littlestone trees).  The VC dimension (\Cref{ch:pac})
determines \emph{whether} learning is possible; the Littlestone dimension
and VCL trees determine \emph{how fast}.}
\label{fig:rate-regimes}
\end{figure}

The structural lesson is a hierarchy of three combinatorial conditions,
each finer than the last:

\begin{center}\small
\renewcommand{\arraystretch}{1.3}
\begin{tabular}{llll}
\toprule
\textbf{Condition} & \textbf{Governs} & \textbf{Chapter} & \textbf{Paradigm} \\
\midrule
$\VC(\mathcal{H}) < \infty$ & Learnability (yes/no) & \Cref{ch:pac} & PAC \\
$\Ldim(\mathcal{H}) < \infty$ & Rate: $e^{-n}$ vs.\ $1/n$ & \Cref{ch:online} & Online \\
No infinite VCL tree & Rate: $1/n$ vs.\ arb.\ slow & This chapter & Universal \\
\bottomrule
\end{tabular}
\end{center}

Each row refines the previous one.  The VC dimension draws the coarsest
line (learnable vs.\ not).  The Littlestone dimension draws a finer line
within the learnable classes (fast vs.\ slow).  The VCL tree draws the
finest line (uniformly slow vs.\ hopelessly slow).

This hierarchy was not predicted by any single tradition.  It emerged from
the trichotomy theorem's proof, which required techniques from both the PAC
and online learning literatures.  The result is a unified picture in which
the three paradigms (PAC, online, and universal) are not separate theories
with accidental parallels, but different projections of a single
combinatorial landscape.


\subsection{Where attention sits in the trichotomy}
\label{sec:attention-trichotomy}

The trichotomy classifies every fixed-architecture attention class, and the classification turns on a single question carried over from \Cref{ch:online}: is the Littlestone dimension of the attention class finite?

\begin{prop}[A fixed attention class is not arbitrarily slow]
\label{prop:attention-trichotomy}
A finite-head attention class on $\R^d$ has finite VC dimension (\cref{prop:attention-pac}), so it has no infinite VCL tree and is not in the arbitrarily-slow regime~\ref{T:arb}: it is universally learnable at a uniform rate.\formal{FLT_Proofs.Criterion.Extended.UniversalLearnable} Its regime is exponential~\ref{T:exp} when its Littlestone dimension is finite and linear~\ref{T:linear} when its Littlestone dimension is infinite.
\end{prop}

\begin{proof}
Finite VC dimension rules out infinite VCL trees (an infinite VCL tree shatters sets of every size, forcing infinite VC dimension), so by \cref{thm:trichotomy} the class lies in regime~\ref{T:exp} or~\ref{T:linear}.  Which of the two is decided by the Littlestone dimension: finite gives the exponential rate via the version-space strategy of \cref{sec:exp-case}, infinite gives the linear rate via the one-inclusion predictor of \cref{sec:linear-case}.
\end{proof}

The rate at which a transformer head learns its target (exponential or merely linear) is therefore governed by exactly the invariant whose finiteness for attention is open (\cref{op:ldim-attention} of \Cref{ch:online}).  A fixed attention class behaves like halfspaces (exponential) if its score-induced order structure has bounded Littlestone depth, and like thresholds on $\R$ (linear) if the temperature lets the router resolve an unbounded ordering; deciding which is the content of \cref{op:attention-regime}.

\section{What This Chapter Established}
\label{sec:ch9-summary}

Two results, one structural lesson.

\begin{enumerate}[leftmargin=*,itemsep=3pt]
\item \textbf{Universal learning rates.}  The optimal learning rate asks a
  harder question than PAC sample complexity: not how many samples for a
  fixed accuracy, but how fast the risk decreases as a function of $n$,
  uniformly over all realizable distributions.

\item \textbf{The trichotomy.}  Every hypothesis class with
  $|\mathcal{H}| \geq 3$ falls into exactly one of three regimes:
  exponential ($e^{-n}$), linear ($1/n$), or arbitrarily slow.  The
  boundaries are determined by the Littlestone dimension and VCL trees.
  There are no intermediate rates.

\item \textbf{The cross-paradigm bridge.}  The Littlestone dimension, born
  in the adversarial online setting of \Cref{ch:online}, controls the
  boundary between exponential and polynomial convergence in the i.i.d.\
  statistical setting.  This was the central open question resolved by
  Bousquet et al.~\cite{BousquetHannekeMoranVanHandelYehudayoff2021}: not
  just what the rate regimes are, but what combinatorial objects govern
  them; and the answer came from an unexpected paradigm.
\end{enumerate}

\section{Open Problems}
\label{sec:ch9-open}

Each problem is anchored to a result above and tagged: a \emph{known unknown} (\textsc{ku}) is an articulated open question; an \emph{unknown known} (\textsc{uk}) is a result in the literature or the verified development not yet absorbed into a clean statement.  Verified handles are named where they exist.

\begin{openprob}[Formalizing the trichotomy, \textsc{uk}]
\label{op:trichotomy-formal}
The regimes are typed (\leanname{FLT_Proofs.Criterion.Extended.UniversalLearnable}, the ordinal dimensions) but the trichotomy theorem is not machine-checked (\cref{rmk:trichotomy-formal}).  Formalize the Bousquet--Hanneke--Moran--Ramanan middle branch (infinite Littlestone, finite VC, hence linear rate) via the doubling/aggregation construction, and complete the boosting lemma in the universal-to-PAC direction (\leanname{FLT_Proofs.Theorem.Separation.universal_imp_pac}) that currently rests on an unfinished step.
\end{openprob}

\begin{openprob}[The trichotomy regime of attention, \textsc{uk}]
\label{op:attention-regime}
\Cref{prop:attention-trichotomy} places a fixed-head attention class in the exponential or linear regime according to its Littlestone dimension.  Decide which: compute the Littlestone dimension of the $k$-head attention class as a function of architecture and temperature, settling whether a transformer head learns its target exponentially or only linearly (tied to \cref{op:ldim-attention} of \Cref{ch:online}).
\end{openprob}

\begin{openprob}[A single dimension for the linear/slow boundary, \textsc{ku}]
\label{op:vcl-characterization}
The boundary between linear and arbitrarily-slow rates is ``no infinite VCL tree''.  Give a clean combinatorial dimension (an ordinal-valued measure on the class) whose finiteness is equivalent to the absence of an infinite VCL tree, the way $\Ldim < \infty$ captures the exponential/linear boundary, building on the typed ordinal dimensions (\leanname{FLT_Proofs.Complexity.Ordinal.OrdinalLittlestoneDim}).
\end{openprob}

\begin{openprob}[Intermediate rates beyond the realizable case, \textsc{ku}]
\label{op:intermediate-rates}
The trichotomy is sharp in the realizable setting: only $e^{-n}$, $1/n$, and arbitrarily slow occur.  Determine the rate landscape in the agnostic and bounded-noise settings: do intermediate rates such as $1/\sqrt{n}$ appear, and is there a corresponding finer trichotomy or a genuine continuum?
\end{openprob}

\begin{openprob}[The rate constants, \textsc{ku}]
\label{op:rate-constants}
\Cref{sec:exp-case,sec:linear-case} give $\Theta(e^{-n})$ and $\Theta(1/n)$ with unspecified constants.  Determine the leading constant in each regime as a function of the Littlestone dimension (exponential) and of the one-inclusion graph degree (linear), matching upper and lower bounds.
\end{openprob}

\begin{openprob}[The multiclass trichotomy, \textsc{uk}]
\label{op:universal-multiclass}
The trichotomy is stated for binary classification.  Determine whether it survives for multiclass and real-valued prediction, with the DS or fat-shattering tree in place of the VCL tree, and whether the three-regime structure persists or splits further (\Cref{ch:extensions}).
\end{openprob}

\begin{openprob}[Distribution-dependent universal rates, \textsc{ku}]
\label{op:adaptive-rates}
The universal rate is uniform over all realizable distributions.  Characterize the finer, distribution-dependent rates (the fastest rate achievable on each fixed $D$) and determine how the per-distribution rate relates to the worst-case regime of \cref{thm:trichotomy}.
\end{openprob}

\begin{openprob}[Temperature and the rate transition, \textsc{uk}]
\label{op:attention-temperature-rate}
\Cref{prop:attention-trichotomy} ties an attention class's regime to its Littlestone dimension.  Determine how the score temperature drives the transition: as the temperature grows and the router approaches a hard argmax, does the Littlestone dimension diverge (pushing the class from exponential to linear), and is there a sharp temperature threshold?
\end{openprob}

\begin{openprob}[The one-inclusion predictor for attention, \textsc{uk}]
\label{op:oneinclusion-attention}
The linear regime is achieved by the one-inclusion predictor (\cref{def:oip}).  Give an efficient implementation of the one-inclusion predictor for a finite-head attention class and bound the constant $C$ in its $C/n$ rate in terms of the architecture, since the abstract predictor is not obviously tractable on attention hypothesis classes.
\end{openprob}

\begin{openprob}[Why online governs statistical rates, \textsc{ku}]
\label{op:universal-online-bridge}
The trichotomy's central surprise is that the Littlestone dimension, an adversarial, online quantity, draws the exponential/linear boundary in the i.i.d.\ setting (\cref{sec:surprise}).  Find a single structural principle, beyond the trichotomy proof, that explains the cross-paradigm bridge: a reason the same combinatorial object should control both worst-case mistakes and average-case rates.
\end{openprob}


\subsection*{Counterfactual exercises: generalizing Figure~\ref{fig:trichotomy-tree}}

Each exercise perturbs the decision tree.  Argue, in the figure's grammar (a two-question ladder of combinatorial tree-conditions of increasing permissiveness, finite Littlestone dimension inside no-infinite-VCL-tree inside all classes, whose two decision nodes are drawn solid where a verified ordinal dimension or the VCL tree definition certifies the threshold, and whose three rate leaves are drawn dashed where the rate value and the partition are Bousquet--Hanneke--Moran--van Handel--Yehudayoff's theorem, cited not proved), that the perturbed picture subsumes the concrete one as an \emph{instance} (a class routed to its regime), a \emph{boundary} (one tree-condition toggled so the regime jumps), or a \emph{frontier} (a dashed, open or literature-only edge).

\begin{enumerate}[leftmargin=*,itemsep=3pt]
\item \textbf{The first diamond is the online-learnable boundary.} Read decision one not as a yes/no on trees but as a verified dimensional threshold. Show this is the figure's defining instance on the exponential branch: ``no infinite Littlestone tree'' is exactly $\Ldim(\mathcal{H}) < \infty$, which is exactly online learnability by the Littlestone characterization (\cref{thm:littlestone-characterization}, \leanname{littlestone_characterization}), so the \textsc{no} edge into the $\Theta(e^{-n})$ leaf is the only rate-regime boundary the kernel certifies; the ordinal refinement used for the universal setting is \leanname{FLT_Proofs.Complexity.Ordinal.OrdinalLittlestoneDim}.

\item \textbf{The VCL tree generalizes the Littlestone tree by shattering.} Replace the second diamond's ``per-path consistency'' by ``per-path shattering''. Show this is the instance that names the second object: a VCL tree (\cref{def:vcl-tree}) demands that the instances along each root-to-leaf path be shattered, not merely consistent with a single hypothesis, so every Littlestone tree is a VCL tree and the converse fails (\cref{rmk:vcl-vs-littlestone}); this is why an infinite VCL tree can exist with $\Ldim < \infty$, making decision two strictly weaker than decision one and the ladder $\{\Ldim<\infty\} \subsetneq \{\text{no }\infty\text{-VCL}\} \subsetneq \{\text{all}\}$ strict at both steps.

\item \textbf{The leaves are typed regimes, the rates are not theorems yet.} Ask the kernel for the statement ``this class learns at rate $\Theta(e^{-n})$''. Show this is the figure's principal frontier: the three regimes are typed (the universal-rate criterion \cref{def:universal-rate} is \leanname{FLT_Proofs.Criterion.Extended.UniversalLearnable}, the ordinal dimensions exist), so the leaf \emph{boxes} are well-typed objects, but the rate \emph{values} $\Theta(e^{-n})$, $\Theta(1/n)$, and arbitrarily-slow, together with the ``exactly one of three'' claim, are not machine-checked; closing them, by formalizing the middle-branch doubling construction and completing the boosting lemma in the universal-to-PAC direction, is exactly \cref{op:trichotomy-formal} (\cite{BousquetHannekeMoranVanHandelYehudayoff2021}), so those edges are dashed.

\item \textbf{Collapse the second diamond to a single dimension.} Suppose someone exhibits an ordinal-valued measure whose finiteness equals ``no infinite VCL tree''. Show this would redraw decision two as a clean threshold matching decision one: the linear/arbitrarily-slow boundary would read ``$\dim < \infty$'' the way the exponential/linear boundary reads ``$\Ldim < \infty$'', unifying both diamonds into one ordinal invariant with three value-ranges; this single-invariant question is open (\cref{op:vcl-characterization}, building on \leanname{FLT_Proofs.Complexity.Ordinal.OrdinalLittlestoneDim}).

\item \textbf{Shrink the class to a single threshold family.} Instantiate $\mathcal{H}$ as thresholds on $\mathbb{N}$. Show this is the boundary that lands on the middle leaf: thresholds have $\VC = 1 < \infty$ (\cref{def:vc-dim}) yet $\Ldim = \infty$, so decision one answers \textsc{yes} and decision two answers \textsc{no} (no infinite VCL tree), placing the class in the linear regime $\Theta(1/n)$; the witness is the threshold class of \Cref{ch:online}, whose $\Ldim=\infty$ is constructive while its $\Theta(1/n)$ rate is the dashed literature claim (\cref{op:trichotomy-formal}).

\item \textbf{Force the bottom leaf with an infinite VC dimension.} Take a class that shatters arbitrarily large finite sets. Show this is the boundary into the arbitrarily-slow regime: infinite VC dimension forces an infinite VCL tree (a tree whose every path shatters its instances), so decision two answers \textsc{yes} and no learner achieves any uniform rate; the leaf is dashed because the lower bound ``for every $R(n)\to 0$, some realizable $D$ defeats every learner'' is the theorem, not kernel-proved (\cref{op:trichotomy-formal}, \cite{BousquetHannekeMoranVanHandelYehudayoff2021}).

\item \textbf{A fixed attention head is never arbitrarily slow.} Replace $\mathcal{H}$ by a finite-head attention class on $\mathbb{R}^d$. Show this is the instance that prunes the bottom branch: finite VC dimension rules out an infinite VCL tree (\cref{prop:attention-trichotomy}, \leanname{FLT_Proofs.Criterion.Extended.UniversalLearnable}), so decision two answers \textsc{no} and the head is universally learnable; whether it sits in the exponential or the linear leaf is decided by its Littlestone dimension, whose finiteness for attention is itself open (\cref{op:attention-regime}, \cref{op:ldim-attention}).

\item \textbf{The Littlestone dimension is an online quantity drawing a statistical line.} Read decision one across paradigms. Show this is the frontier that makes the tree surprising: the same invariant that bounds worst-case mistakes in the adversarial online setting of \Cref{ch:online} draws the exponential/linear boundary in the i.i.d.\ statistical setting (\cref{sec:surprise}), and a single structural reason for this cross-paradigm coincidence, beyond the trichotomy proof, is open (\cref{op:universal-online-bridge}).

\item \textbf{Leave the realizable case and the partition may refine.} Relabel the leaves for agnostic or bounded-noise data. Show this is the frontier that may add regimes: the trichotomy is sharp only in the realizable setting, and whether intermediate rates such as $1/\sqrt{n}$ appear, splitting a leaf or introducing a fourth, is open (\cref{op:intermediate-rates}); the figure's ``no fourth regime'' annotation is therefore a realizable-case claim, not a universal one.

\item \textbf{Swap the VCL tree for a multiclass tree.} Replace binary labels by a finite or real-valued label set and the VCL tree by the DS or fat-shattering tree. Show this is the frontier that tests the ladder's stability: whether the three-regime structure persists with the new combinatorial object in the second diamond, or splits further, is open (\cref{op:universal-multiclass}, \Cref{ch:extensions}); decision two's object is paradigm-relative, while decision one's $\Ldim$ threshold is the fixed verified rung.
\end{enumerate}
